\section{Einleitung}
Das Bohr-van-Leeuwen Theorem besagt, dass bei der Anwendung der klassischen Statistik die Magnetisierung
in eines Festkörpers im thermischen Gleichgewicht null ist. Das führt unabdingbar zu der Schlussfolgerung, dass
magnetische Phänomene wie Ferromagnetismus, Diamagnetismus und Paramagnetismus rein quantenmechanisch
erklärt werden können. \\
Ferromagnetische Materialien nehmen dabei eine besondere Stellung ein,
da sie im Gegensatz zu Dia- und Paramagnetika auch ohne äußeres
Magnetfeld eine spontane Magnetisierung besitzen. Diese entsteht durch
die parallele Ausrichtung mikroskopischer magnetischer Momente in
sogenannten Weissschen Bezirken infolge der quantenmechanischen
Austauschwechselwirkung. Charakteristisch für Ferromagneten sind
insbesondere Hystereseeffekte, Remanenz und Koerzitivfeldstärke,
welche sie für technische Anwendungen wie Transformatoren,
Elektromotoren oder magnetische Datenspeicherung besonders relevant
machen.\\
Trotz der rein quantenmechanischen Ursache des Magnetismus lassen sich in der makroskopischen Welt
verschiedene magnetische Eigenschaften von Materialien beobachten. 
Die in diesem Versuch untersuchten Materialien gehören zur Klasse
der technisch relevanten Ferromagnetika und unterscheiden sich
insbesondere in ihrer Hysteresekurve sowie in ihren magnetischen
Kenngrößen. Während bei schwach magnetisierbaren Stoffen typischerweise
die magnetische Suszeptibilität oder die relative Permeabilität als
Materialkonstanten betrachtet werden, stehen bei Ferromagneten aufgrund
der nichtlinearen und hysteresebehafteten Magnetisierungskurve andere
Größen im Vordergrund.
Die im Experiment bestimmten magnetischen Kenngrößen wie die Remanenzflussdichte
\( B_r \), die Koerzitivfeldstärke \( H_c \) sowie die Fläche der
Hystereseschleife sind streng genommen keine Materialkonstanten im
thermodynamischen Sinn, da sie von der Probenform, der
Vormagnetisierung und den Messbedingungen abhängen.